\begin{song}{Es War An Einem Sommertag}{Arno Clauss}

    \begin{verse}  
        Es [C]war an [G]einem [C]Sommer[G]tag, [Am]irgendwann und [E]irgend[Am]wo. \\
        Da [C]tönte [G]plötzlich [C]Trommel[G]schlag, und [Am]Flöten[E]spiel klang [Am]froh. \\
        Es [G]war eine muntre, [C]bunte Schar, die [G]machte vor dem Rathaus [C]Halt \\
        Den [G]Grund, weshalb sie ge[C]kommen war, [Am]erfuhren die [E]Leute [Am]bald.
    \end{verse}  

    \begin{chorus}
        (pfeifen)
        % [C]e-c-d-e-[G]d--g--[C]e-c-d-e-[G]d--h--[Am]c-a-h-c-[E]h--c-[Am]a----
    \end{chorus}

    \begin{verse}  
        Ein [C]Mann mit [G]einem [C]Feder[G]hut rief: [Am]Männer! [E]Hört mir [Am]zu! \\
        [C]Ich ver[G]sprech' euch [C]Geld und [G]Gut und [Am]Ehre [E]noch da[Am]zu! \\
        Der [G]Kaiser braucht euch, [C]reiht euch ein! [G]Hängt nicht an Weib und [C]Haus! \\
        Es [G]wird auch gar nicht [C]lange sein – zieht [Am]mit ins [E]Feld hi[Am]naus!
    \end{verse}  

    \refchorus

    \begin{verse}  
        Im [C]Wirtshaus [G]war das [C]Trinken [G]frei, be[Am]zahlt von des [E]Kaisers [Am]Gold. \\
        Und [C]während [G]dieser [C]Zeche[G]rei trat [Am]mancher in des [E]Kaisers [Am]Sold. \\
        Gab [G]seiner Frau den [C]Abschiedskuss, [G]versuchte als Soldat sein [C]Glück, \\
        sah [G]nicht des Werbers [C]Pferdefuß und [Am]kehrt' nie [E]mehr zu[Am]rück.
    \end{verse}  

    \refchorus

    \begin{verse}  
        Mit [C]Flöten[G]spiel und [C]Trommel[G]schlag ging’s [Am]früh am [E]Morgen [Am]fort. \\
        Die [C]Schar ward [G]größer, [C]denn es [G]lag am [Am]Weg noch so [E]mancher [Am]Ort. \\
        Der [G]Werber mit dem [C]Federhut macht [G]sein Geschäft nicht [C]schlecht, \\
        ver[G]sprach noch vielen [C]Geld und Gut – dem [Am]Kaiser, [E]dem war’s [Am]recht.
    \end{verse}  

    \refchorus

    \begin{verse}  
        Die [C]Jahre [G]gingen [C]in das [G]Land, und [Am]von der [E]großen [Am]Schar \\
        gab’s [C]keine:n, [G]der nach [C]Hause [G]fand, wie [Am]er ge[E]gangen [Am]war: \\
        Die [G]eine ließ ihr [C]Bein im Feld, [G]blind kam ein andrer [C]an. \\
        Die [G]meisten hatte der [C]Tod gefällt, der [Am]jede [E]Schlacht ge[Am]wann.
    \end{verse}  

    \refchorus

    \begin{verse}  
        Die [C]letzten [G]Tränen [C]waren [G]kaum ge[Am]weint, da [E]waren [Am]sie \\
        auch [C]schon ver[G]gessen [C]wie ein [G]Traum, die [Am]Menschen [E]lernen [Am]nie. \\
        Und [G]dann an einem [C]Sommertag, irgend[G]wann und irgend[C]wo, \\
        da [G]tönte plötzlich [C]Trommelschlag, und [Am]Flöten[E]spiel klang [Am]froh.
    \end{verse}  

    \refchorus

\end{song}
